\section{Implementation and reproducibility}
\label{app:implementation}

\subsection{Current implementation}
The source, complete numerical evidence and compiled manuscript are at
\begin{center}
\url{https://github.com/zhixin0825/spectral-lrt}.
\end{center}
Algorithm~1 is implemented by \texttt{SpectralLikelihood.decode()} in
\begin{center}
\small\path{experiments/20261010_general_density/spectral_likelihood.py}.
\end{center}
It accepts an $n\times K$ eigenvector array and a length-$K$ array of the corresponding signed eigenvalues. The decoder receives neither the adjacency matrix, true labels, true block probabilities nor node degrees. Retaining the top-$|\lambda|$ eigenpairs is the preceding spectral step.

The default \texttt{score\_family='bernoulli'} evaluates \eqref{eq:working-score}. The floor coefficient defaults to $c_0=10^{-4}$; its theoretical guarantee requires the explicit admissibility condition in Section~\ref{sec:algorithm}, not merely positive block probabilities. The spectral scale is $\max\{1,\max_a|\lambda_a|\}$, and the reconstructed entries are clipped on both sides. The original adjacency is not regularized. The optional \texttt{score\_family='sparse'} is a score control on precisely the same $q$, and is not the default general-density algorithm.

Growing uses at most 100 exact M-steps per stage. A fit stops after at least one M-step if hard labels are unchanged and its absolute soft-loss change is at most $10^{-9}n$. Replacement uses at most $\lceil\log n\rceil$ rounds, tries every component deletion against the same incumbent, initializes each trial's weights uniformly, and uses one M-step per trial by default. A full round with no improvement stops the search. This rule is justified by the contraction certificate in Appendix~\ref{app:replacement-proof}. Calling \texttt{decode()} always executes the final E-step and one profile/weight M-step before returning labels.

The implementation uses the unconstrained simplex weight update. Zero-mass components retain their previous profile. Log-sum-exp arithmetic evaluates responsibilities and the soft divergence loss $F$ without exponentiating very negative raw log-likelihoods. Theoretical divergences are nonnegative; the code clips negative floating-point residue at zero. Weighted means, including the initial global mean, are clipped back to the same feasible interval solely to correct roundoff in convex combinations. The theorem is an exact-arithmetic result; recorded numerical checks assess implementation accuracy to floating-point precision.

The user may increase either fitting budget. Explicitly requesting fewer replacement rounds is an experimental option and does not automatically retain the stated general-$K$ initialization guarantee. The proven default budget is $\lceil\log n\rceil$, with valid early stopping after a full unsuccessful round.

\subsection{Current finite-graph and population evidence}
The complete 60-graph design is specified by
\begin{center}
\small\path{experiments/20261010_general_density/benchmark_results/protocol.json}.
\end{center}
It fixes all graph seeds, model families, density settings, method parameters, and hashes of the code used at fitting time. Each graph directory retains the original $A$, $P$, true labels for evaluation, exact retained $(U,\Lambda)$, all fitted profiles and weights, every candidate's hard loss, and full EM objective traces. Ground truth is not used for fitting, selecting replacements, or stopping. The aggregate CSV contains all 360 graph--method combinations.

The archived fitting engine is retained as \texttt{engine\_at\_fitting.py}. A subsequent boundary fix clips the rounded initial mean when a graph has all reconstructed entries at the floor. The maintenance record checks that this expression is bitwise unchanged on every archived graph, so the saved benchmark is unaffected. The original fitting hashes remain in the frozen protocol; maintenance provenance is recorded separately.

The five-community population calculation and the exact rational interval certificate are in
\begin{center}
\small\path{experiments/20261010_general_density/population_audit/}.
\end{center}
They record the positive-definiteness certificate, strict all-node-gain comparisons, the limiting one-update-per-stage trajectory, and replacement trial results. The distinction between a population limiting path, a finite-graph theorem, and fully iterated growing is maintained in the report. The manifest \texttt{SHA256SUMS} identifies the saved numerical records.

\subsection{Historical experiments and entry points}
The earlier 240-graph comparison is preserved at
\path{experiments/20261009_residual_likelihood_seeding/}.
The \texttt{residual\_seed.py} module retains several sparse-score methods, using
\[
 \ell_i^{\mathrm{sparse}}(\widehat p)=\sum_jq_{ij}\log\widehat p_j-\sum_j\widehat p_j.
\]
Its current default \path{method='global_gain_certified'} calls the separate sparse threshold-certificate implementation \texttt{Work.iterative\_certified\_growing()}. It uses an adaptive lower-only floor, a uniform-weight refit of the current centers as well as $K$ replacement trials, acceptance threshold $(\sum_{ij}q_{ij})/\sqrt{\log n}$, and a cap of $\lceil(\log n)^2\rceil$ accepted rounds. Its general-$K$ proof for the sparse range $p\to0$ with $p=\Omega(\log n/n)$ and its 24-graph implementation checks are retained as a companion version. They preceded the complete Bernoulli density extension and direct soft-loss contraction presented here, which allow the full range \eqref{eq:density-condition}.

In that module, \path{method='global_gain'} is the growing-only branch, \path{method='global_gain_restarts'} selects independent residual starts, and \path{method='repair'} is the historical single-pass repair. None identifies the current full-Bernoulli Algorithm~1, whose entry point is specified above. The actual 240-graph runs used the earlier lower floor $10^{-4}\log n/n$ and their recorded protocol, not the adaptive-floor default now present in the companion module. Their profiles were not constrained to lie below one. The preserved benchmark runner explicitly selects \texttt{floor\_mode='historical'} for these original arms; their earlier source is also accessible in repository history.

The archived 240-graph fits use an exact weight M-step with floor $f=10^{-8}$. For a fixed $0\le f<1/K$, it solves
\[
 \max_{\pi_a\ge f,\,\sum_a\pi_a=1}\sum_aN_a\log\pi_a,
 \qquad \pi_a=\max\{f,N_a/\lambda\},
 \quad\sum_a\pi_a=1.
\]
An active-set solver determines $\lambda$; naive clipping and renormalization is not this optimizer. The current main theorem and implementation use $f=0$.

Historical fits stop only when labels are unchanged, relative profile and absolute weight changes are at most $10^{-3}$, and the objective increment is at most $10^{-6}(1+|\mathcal L-G|)$, where $G=\sum_{ij}\log\Gamma(q_{ij}+1)$ is an observation-only tolerance constant. The cap is 300 updates. Historical top-subset and peeling controls, their subset sizes and stopping rules are defined by the archived code. They are not steps of Algorithm~1.

\subsection{Build and validation}
The active manuscript contains five main sections, statistical and deterministic proof appendices, implementation details, and references. PDFLaTeX, BibTeX and two further PDFLaTeX passes resolve the bibliography and cross references. Build provenance records the source revision and source/PDF checksums. The manuscript is rendered for a page-by-page layout check. These syntax, reference and layout checks are separate from the mathematical audit and numerical evidence described above.
